\documentclass[11pt,a4paper]{article}

\usepackage[T1]{fontenc}
\usepackage[utf8]{inputenc}
\usepackage{lmodern}
\usepackage{microtype}
\usepackage[margin=2.5cm]{geometry}
\usepackage{amsmath,amssymb}
\usepackage{booktabs}
\usepackage{array}
\usepackage{enumitem}
\usepackage{fancyvrb}
\usepackage{listings}
\usepackage{upquote}
\usepackage{forest}
\usepackage[hidelinks]{hyperref}

\setlength{\parindent}{0pt}
\setlength{\parskip}{0.55em}

% ---------- code ----------
\lstset{
  language=Java,
  basicstyle=\ttfamily\small,
  keywordstyle=\bfseries,
  commentstyle=\itshape,
  columns=fullflexible,
  keepspaces=true,
  upquote=true,
  showstringspaces=false,
  frame=lines,
  framesep=4pt,
  xleftmargin=0pt,
  xrightmargin=0pt,
  aboveskip=0.8em,
  belowskip=0.8em,
  numbers=none
}
\newcommand{\code}[1]{\texttt{#1}}
\fvset{fontsize=\small}

% ---------- checkpoints (renders of the world) ----------
\newenvironment{checkpoint}[1]
  {\begin{tabular}[t]{@{}l@{}}\footnotesize #1\\[2pt]}
  {\end{tabular}}

% ---------- questions ----------
\newcounter{qcount}
\newcommand{\Q}{\refstepcounter{qcount}\textbf{Q\theqcount.}\ }
\newcommand{\task}[1]{\par\medskip\textbf{Task #1.}\ }
\newcommand{\beyond}{\textsuperscript{$\star$}}

\begin{document}

\begin{center}
  {\Large\bfseries Mini-project: a falling-sand sandbox}\\[0.6em]
  Introduction to Programming --- Inheritance\\[0.3em]
  {\small Estimated time: 4 hours \quad$\cdot$\quad}
\end{center}

\vspace{0.5em}

\section*{Overview}

You will build a small \emph{falling-sand} simulator: a rectangular world of cells in which sand
piles up, water spreads, oil floats on water, wood catches fire and smoke rises and fades away.
The world is printed in the terminal as ASCII art, one frame per tick, so you can watch it evolve.

The interesting part is not the physics, which is deliberately naive, but the \emph{shape of the
code}. Each kind of matter is a class in an inheritance hierarchy. The world only ever manipulates
variables of type \code{Element}, and every difference in behaviour comes from redefining methods
(dynamic binding). The project exercises, in order of appearance:
\begin{itemize}[nosep]
  \item the \emph{is-a} relation and the definition of sub-classes with \code{extends} (lesson 10);
  \item access rights \code{private}, \code{protected}, \code{public} in a hierarchy (lesson 11);
  \item redefinition of methods, \code{super.}\emph{method}, and attribute hiding, i.e.\ \emph{masquage} (lesson 12);
  \item constructors in a hierarchy, \code{super(\ldots)}, \code{this(\ldots)}, the order of construction (lesson 13);
  \item static versus dynamic binding and heterogeneous collections (lesson 14);
  \item and, from earlier weeks: encapsulation, copy constructors, \code{toString}, reference
        semantics, two-dimensional arrays and strings.
\end{itemize}

\paragraph{Time budget.}
\begin{center}
\begin{tabular}{@{}llr@{}}
\toprule
Part & Content & Time \\
\midrule
A & Warm-up on paper: predict, then justify & 30 min \\
B & Elements and the world & 35 min \\
C & Falling particles & 40 min \\
D & Liquids & 40 min \\
E & Fire and smoke & 45 min \\
F & Parallel futures: copying a world & 30 min \\
G & Your own scene, design table, reflection & 20 min \\
\midrule
 & Total & 4 h \\
\bottomrule
\end{tabular}
\end{center}

\paragraph{How to check yourself.}
Every rule in Section~\ref{sec:spec} is deterministic: there is no randomness. The expected
renders shown in each part (the \emph{checkpoints}) were produced by a reference implementation of
exactly this specification. If your output differs, find the \emph{first} tick at which it differs:
that tick and that cell are your best debugging clue. Questions are numbered
Q1, Q2, \ldots; answer them in a file \code{answers.txt} \emph{before} looking anything up.
A star\beyond{} marks a question that goes slightly beyond what the slides state explicitly and
asks you to reason from the rules you know.

\paragraph{Design constraints.} These are part of the specification and are what the project
actually tests. Re-read them at the start of every part.
\begin{enumerate}[label=\textbf{R\arabic*.},leftmargin=2.6em]
  \item \textbf{Course material only.} Use only BOOC I and weeks 1--3 of BOOC II. In particular:
        no \code{abstract}, no \code{interface}, no \code{instanceof}, no casts between class
        types, no \code{getClass()}, no \code{static} attributes, no exceptions (except in the
        helper of Appendix~\ref{app:helper}). Static methods are allowed only in your main class,
        as in BOOC I. If you feel that one of these would make the code better, write a comment
        saying where and why: this is a valuable observation for weeks 4--6.
  \item \textbf{Encapsulation.} All attributes are \code{private}. If you declare any member
        \code{protected}, justify it in a one-line comment next to it.
  \item \textbf{No type tests.} \code{World}, \code{Particle}, \code{Liquid}, \code{Gas} and
        \code{Fire} never test what kind of element they are dealing with: no comparison of
        \code{getName()} or \code{getSymbol()} to decide behaviour (the only exceptions are
        \code{World.count} and the parsing method). Differences in behaviour must come from
        redefinition.
  \item \textbf{Concrete classes are named in few places.} The names \code{Air}, \code{Wall},
        \code{Wood}, \code{Fire}, \code{Sand}, \code{Water}, \code{Oil}, \code{Smoke} may appear only
        in: the parsing method of \code{World}; \code{Fire} (which creates \code{Fire} and
        \code{Smoke}); \code{Gas} (which creates \code{Air}); each class's own constructors and
        \code{copy()} method; and your main class.
  \item \textbf{No duplication.} The falling algorithm of Section~\ref{sec:particles} is written
        exactly once. Each \code{toString} that extends another one uses \code{super.toString()}.
  \item \textbf{Constructors.} \code{Element} has no default constructor. Every constructor of
        every sub-class begins with an explicit \code{super(\ldots)} or \code{this(\ldots)}.
\end{enumerate}

% =====================================================================
\section{Part A --- Warm-up on paper (30 min)}

Answer without a computer, then check by compiling. For every answer, write \emph{one sentence of
justification} naming the rule you used (``static type'', ``dynamic type'', ``implicit
\code{super()}'', \ldots). A correct output with a wrong reason counts as wrong.

\subsection*{A1. Order of construction}
\begin{lstlisting}
class A {
    public A()      { System.out.print("A"); }
    public A(int x) { this(); System.out.print("a" + x); }
}
class B extends A {
    public B()      { System.out.print("B"); }
    public B(int x) { super(x + 1); System.out.print("b" + x); }
}
class C extends B {
    public C()      { super(2); System.out.print("C"); }
}
\end{lstlisting}
\Q What do \code{new C()} and \code{new B()} print?

\Q The constructor \code{A()} is deleted and nothing else changes. Which lines no longer compile,
and why? Which constructors are unaffected?

\subsection*{A2. A band of instruments}
\begin{lstlisting}
class Instrument {
    private String name;
    public Instrument(String name) { this.name = name; }
    public String getName()  { return name; }
    public String describe() { return "instrument " + name; }
    public String play()     { return "..."; }
    public String concert()  { return describe() + " plays " + play(); }
}
class StringInstrument extends Instrument {
    private int strings;
    public StringInstrument(String name, int strings) {
        super(name);
        this.strings = strings;
    }
    public String describe() {
        return super.describe() + " with " + strings + " strings";
    }
}
class Guitar extends StringInstrument {
    public Guitar() { super("guitar", 6); }
    public String play() { return "strum"; }
}
class Harp extends StringInstrument {
    public Harp() { super("harp", 47); }
    public String describe() { return "a big " + super.describe(); }
    public String play()     { return "pling" + super.play(); }
}
class Drum extends Instrument {
    public Drum() { super("drum"); }
    public String play() { return "boom"; }
}
// in main:
Instrument[] band = { new Guitar(), new Harp(), new Drum() };
for (int i = 0; i < band.length; ++i) {
    System.out.println(band[i].concert());
}
\end{lstlisting}
\Q Give the three printed lines. For the second line, list in order every method body that is
executed, naming its class.

\Q \code{concert()} is written once, in \code{Instrument}, a class that knows nothing about harps.
Explain precisely why it nevertheless calls \code{Harp}'s \code{describe()}.

\Q Does the statement \code{StringInstrument s = band[0];} compile? \code{band[0]} is a guitar,
after all.

\Q The constructor of \code{Drum} is deleted. Does the program still compile? What if, in
addition, \code{Instrument} received a second constructor \code{Instrument()}?

\subsection*{A3. Hiding an attribute\beyond}
\begin{lstlisting}
class P {
    protected int x = 1;
    public int getX() { return x; }
}
class Q extends P {
    private int x = 2;                       // masquage
    public int getX() { return x; }
    public int sum()  { return x + super.x + super.getX(); }
}
// in main (same package):
P p = new Q();
Q q = new Q();
System.out.println(p.getX());
System.out.println(p.x);
System.out.println(q.sum());
\end{lstlisting}
\Q Predict the three lines. Which of them depend on the static type of \code{p}, and which on its
dynamic type? Use your answer to formulate a rule: ``methods are \ldots, attributes are \ldots''.
Then explain why the course advises against \emph{masquage}.

\subsection*{A4. \texttt{protected} across packages}
Four files, in two packages \code{zoo} and \code{park}.

\noindent
\begin{minipage}[t]{0.48\linewidth}
\begin{lstlisting}
// file zoo/Animal.java
package zoo;
public class Animal {
    protected int energy;
    private int age;
    public int getAge() {
        return age;
    }
}
\end{lstlisting}
\end{minipage}\hfill
\begin{minipage}[t]{0.48\linewidth}
\begin{lstlisting}
// file park/Lion.java
package park;
import zoo.Animal;
public class Lion extends Animal {
    public void eat() {
        energy = energy + 5; // (1)
        age = age + 1;       // (2)
    }
}
\end{lstlisting}
\end{minipage}

\noindent
\begin{minipage}[t]{0.48\linewidth}
\begin{lstlisting}
// file park/Keeper.java
package park;
public class Keeper {
    public void feed(Lion l) {
        l.energy += 1;       // (3)
        int a = l.getAge();  // (4)
    }
}
\end{lstlisting}
\end{minipage}\hfill
\begin{minipage}[t]{0.48\linewidth}
\begin{lstlisting}
// file zoo/Vet.java
package zoo;
import park.Lion;
public class Vet {
    public void check(Lion l) {
        l.energy = 0;        // (5)
    }
}
\end{lstlisting}
\end{minipage}

\Q Which of the lines (1)--(5) compile? Justify each one with the rule from lesson~11.

\subsection*{A5. Modelling}
\Q For each pair, decide whether the right relation is \emph{is-a} (inheritance), \emph{has-a}
(an attribute) or \emph{neither}, and justify in one line:
(\code{World}, \code{Element}); (\code{Water}, \code{Liquid}); (\code{Sand}, \code{Element});
(\code{Fire}, \code{Smoke}); (\code{Liquid}, \code{Powder}); (\code{Wood}, \code{Fire}).
For (\code{Liquid}, \code{Powder}), note that a liquid does fall exactly like a powder before
flowing sideways: is reusing code a sufficient reason to inherit?

% =====================================================================
\section{The world: specification}\label{sec:spec}

\subsection{Grid, ticks and elements}

The world is a grid of \(R\) rows and \(C\) columns. Cell \((r,c)\) has row \(r\in\{0,\dots,R-1\}\)
and column \(c\in\{0,\dots,C-1\}\); row \(0\) is at the \emph{top}. Every cell always contains
exactly one element; empty space is the element \code{Air}, never \code{null}. Positions outside the
grid do not exist: any move or neighbour outside the grid is simply ignored.

Time advances in \emph{ticks}. The world stores the current tick, which is \(0\) after construction
and is incremented by one at the beginning of each step.

\begin{center}
\begin{forest}
  for tree={
    font=\ttfamily\small, draw, inner sep=3pt, s sep=6pt, l sep=14pt,
    edge={latex-}, parent anchor=south, child anchor=north, anchor=north
  }
  [Element
    [Air] [Wall] [Wood] [Fire]
    [Particle, dashed
      [Powder, dashed [Sand]]
      [Liquid, dashed [Water] [Oil]]
      [Gas, dashed [Smoke]]
    ]
  ]
\end{forest}

\smallskip
{\small Arrows point to the super-class. Dashed classes never appear directly in a scene:
they hold behaviour shared by their sub-classes.}
\end{center}

\begin{center}
\small
\begin{tabular}{@{}llcrcllp{2.6cm}@{}}
\toprule
Class & name & symbol & density & fluid & flammable & extinguishes & other \\
\midrule
\code{Air}   & \code{air}   & \verb*| | & 2   & yes & no               & no  & \\
\code{Wall}  & \code{wall}  & \verb|#|  & 100 & no  & no               & no  & \\
\code{Wood}  & \code{wood}  & \verb|=|  & 30  & no  & yes, burn 10 & no  & \\
\code{Fire}  & \code{fire}  & \verb|*|  & 0   & no  & no               & no  & 6 ticks by default \\
\code{Sand}  & \code{sand}  & \verb|.|  & 10  & no  & no               & no  & \\
\code{Water} & \code{water} & \verb|~|  & 6   & yes & no               & yes & spread 3 \\
\code{Oil}   & \code{oil}   & \verb|o|  & 4   & yes & yes, burn 3  & no  & spread 1 \\
\code{Smoke} & \code{smoke} & \verb|"|  & 1   & yes & no               & no  & lifetime 8 \\
\bottomrule
\end{tabular}

\smallskip
{\small The symbol of \code{Air} is a space. Any character not in this table is read as air.}
\end{center}

\subsection{Public interface of \texttt{Element}}

\begin{lstlisting}
public Element(String name, char symbol, int density)
public String  getName()
public char    getSymbol()
public int     getDensity()
public void    markUpdated(int tick)    // "I was updated at this tick"
public boolean wasUpdatedAt(int tick)   // true iff last mark == tick
public boolean isFluid()                // see the table
public boolean isFlammable()
public int     getBurnTime()            // 0 for non-flammable elements
public boolean extinguishesFire()
public boolean update(World w, int row, int col)
public String  toString()
\end{lstlisting}

\code{update} is called by the world when it is this element's turn to act; the element is at
\((\mathit{row},\mathit{col})\) in \code{w}. It returns \code{true} if and only if, during this call,
the element moved to another cell or was replaced. In \code{Element}, it does nothing and returns
\code{false}: this is the general behaviour of matter that never acts (air, walls, wood).

\code{toString()} returns \textit{name}\verb*| '|\textit{symbol}\verb*|' d=|\textit{density},
for instance \verb|wall '#' d=100|. Some sub-classes extend this text (Parts D and E).

\subsection{Public interface of \texttt{World}}

\begin{lstlisting}
public World(String[] lines)            // one string per row, equal lengths
public int     getRows()
public int     getCols()
public int     getTick()
public boolean isInside(int row, int col)
public Element get(int row, int col)    // precondition: isInside(row, col)
public void    set(int row, int col, Element e)
public void    swap(int r1, int c1, int r2, int c2)
public void    step()                   // one tick
public void    run(int n)               // n ticks
public String  renderRow(int row)       // the symbols of one row
public String  render()                 // all rows, each followed by '\n'
public int     count(String name)       // cells whose element has this name
\end{lstlisting}

\paragraph{Rule W1 (marking).} \code{set(r, c, e)} stores \code{e} in the cell \emph{and marks
\code{e} as updated at the current tick}. \code{swap} exchanges two cells and marks nothing. The
constructor fills the grid using \code{set}. Consequence: an element created during tick \(t\)
acts for the first time at tick \(t+1\).

\paragraph{Rule W2 (one step).} \code{step()} does the following.
\begin{enumerate}[nosep]
  \item Increment the tick; call its new value \(t\).
  \item For \(r = R-1\) down to \(0\) (bottom to top), and for \(c = 0\) up to \(C-1\) (left to
        right): let \(e\) be the element in cell \((r,c)\) \emph{at that moment}. If
        \(e\) was updated at tick \(t\), skip it. Otherwise mark \(e\) as updated at \(t\), then
        call \code{e.update(this, r, c)}.
\end{enumerate}
The mark guarantees that an element acts at most once per tick, even if it moves to a cell
that has not been scanned yet.

\subsection{Particles}\label{sec:particles}

A \code{Particle} is an element that moves by swapping places with fluid neighbours. It has
\begin{lstlisting}
public int     getGravity()   // +1 (falls); gases: -1 (rises)
public boolean canEnter(Element target)
public boolean tryMove(World w, int row, int col, int dRow, int dCol)
\end{lstlisting}

\paragraph{Rule P1 (who gives way).} Let \(g\) be the gravity of particle \(p\). Then
\(p\) can enter \(\mathit{target}\) if and only if \(\mathit{target}\) is fluid and
\[
  g > 0 \;\Rightarrow\; d(\mathit{target}) < d(p),
  \qquad
  g < 0 \;\Rightarrow\; d(\mathit{target}) > d(p),
\]
where \(d\) denotes density. \code{tryMove(w, r, c, dr, dc)} checks that \((r+dr, c+dc)\) is inside
the grid and that the particle can enter the element there; if so it swaps the two cells and
returns \code{true}, otherwise it returns \code{false} and changes nothing.

\paragraph{Rule P2 (general movement).} The general \code{update} of a particle at \((r,c)\) tries
the following moves in order and stops at the first one that succeeds:
\[
  (g,\,0), \qquad (g,\,s), \qquad (g,\,-s),
  \qquad\text{where } s = \begin{cases} -1 & \text{if the current tick is even},\\
                                         +1 & \text{if the current tick is odd.}\end{cases}
\]
That is: straight down (or up, for a gas), then diagonally, alternating which diagonal is tried
first so that piles stay symmetric. It returns \code{true} iff a move succeeded.

% =====================================================================
\section{Part B --- Elements and the world (35 min)}

\task{B1} Write \code{Element} with the interface above and private attributes. The mark can be
stored as the last tick at which the element was updated (initially \(-1\)).

\task{B2} Write \code{Air} and \code{Wall}. Each needs only a constructor without parameters and,
for \code{Air}, one redefinition. Look at lesson~13's \code{Carre}: ``and that's all''.

\task{B3} Write \code{World}: a private two-dimensional array of \code{Element}, the tick, and
the interface above. Parsing goes through a private method \code{Element createElement(char symbol)}
(an instance method, called from the constructor). For now, it knows only \verb|#| and the space;
you will extend it as classes appear. Implement \code{step()} completely now, following rule W2:
it does not depend on which elements exist.

\paragraph{Checkpoint B.} Build the scene \code{drop} of Appendix~\ref{app:scenes}. Then
\code{render()} returns the scene unchanged (with \verb|.| read as air for now);
\code{count("wall")} is \(24\); \code{System.out.println(w.get(0, 0))} prints \verb|wall '#' d=100|
and \code{w.get(1, 1)} prints \verb*|air ' ' d=2|. Explain in one line why \code{println} prints
these texts although you never called \code{toString()} yourself.

\Q \code{Element} has no default constructor (R6). What would go wrong, in the context of
lesson~13, if it had one that set the name to \code{"?"}? Think of a sub-class whose author
forgets to call \code{super(\ldots)}.

% =====================================================================
\section{Part C --- Falling particles (40 min)}

\task{C1} Write \code{Particle} according to Section~\ref{sec:particles}, then \code{Powder} and
\code{Sand}. \code{Sand} contains nothing but its constructor. Add \verb|.| to the parser.

\task{C2} In your main class, write a static function
\code{boolean conserves(World w, String name, int n)} that runs \code{w} for \code{n} ticks and
returns whether \code{count(name)} is unchanged. Sand must be conserved in every scene without fire.

\paragraph{Checkpoint C.} Scene \code{drop}, ticks \(0\) to \(4\):
% scene drop
\begin{center}
\begin{checkpoint}{$t=0$}
\begin{BVerbatim}
#######
#  .  #
#  .  #
#  .  #
#     #
#     #
#######
\end{BVerbatim}
\end{checkpoint}
\hspace{1.6em}
\begin{checkpoint}{$t=1$}
\begin{BVerbatim}
#######
#     #
#  .  #
#  .  #
#  .  #
#     #
#######
\end{BVerbatim}
\end{checkpoint}
\hspace{1.6em}
\begin{checkpoint}{$t=2$}
\begin{BVerbatim}
#######
#     #
#     #
#  .  #
#  .  #
#  .  #
#######
\end{BVerbatim}
\end{checkpoint}
\hspace{1.6em}
\begin{checkpoint}{$t=3$}
\begin{BVerbatim}
#######
#     #
#     #
#     #
#  .  #
#  .. #
#######
\end{BVerbatim}
\end{checkpoint}
\hspace{1.6em}
\begin{checkpoint}{$t=4$}
\begin{BVerbatim}
#######
#     #
#     #
#     #
#     #
# ... #
#######
\end{BVerbatim}
\end{checkpoint}
\end{center}

Scene \code{hourglass} at ticks \(5\), \(20\) and \(60\), with \code{count("sand")} \(=24\)
throughout:
% scene hourglass
\begin{center}
\begin{checkpoint}{$t=5$}
\begin{BVerbatim}
###########
#    .....#
## ......##
###.....###
####...####
#####.#####
#### . ####
###  .  ###
##   .   ##
#    .    #
###########
\end{BVerbatim}
\end{checkpoint}
\hspace{1.6em}
\begin{checkpoint}{$t=20$}
\begin{BVerbatim}
###########
#        .#
##      .##
###    .###
####  .####
#####.#####
####.. ####
### ....###
##...... ##
# ....... #
###########
\end{BVerbatim}
\end{checkpoint}
\hspace{1.6em}
\begin{checkpoint}{$t=60$}
\begin{BVerbatim}
###########
#         #
##       ##
###     ###
####   ####
##### #####
####...####
###.....###
##.......##
#.........#
###########
\end{BVerbatim}
\end{checkpoint}
\end{center}


\Q (\emph{What changes if\ldots?}) Suppose rule W2 scanned rows from top to bottom instead.
Predict the scene \code{drop} after one tick, \emph{then} check by temporarily changing your code.
Why is bottom-to-top the natural order for falling matter, and which kind of particle will
prefer the opposite order?

\Q \code{Powder} currently adds nothing to \code{Particle}. Give two reasons, one about modelling
and one about the future of the code, to keep it in the hierarchy.

% =====================================================================
\section{Part D --- Liquids (40 min)}

A \code{Liquid} is a fluid particle with an integer \emph{spread} (given to its constructor) and a
private \emph{direction} \(\delta\in\{-1,+1\}\), initially \(+1\) (to the right).

\paragraph{Rule L1 (liquid movement).} The \code{update} of a liquid at \((r,c)\):
\begin{enumerate}[nosep]
  \item tries the general movement of rule P2; if it succeeds, returns \code{true};
  \item otherwise \emph{flows}: repeatedly moves one cell in direction \(\delta\), i.e.\ performs
        \code{tryMove(w, r,} \(c'\)\code{, 0,} \(\delta\)\code{)} from its \emph{current} column \(c'\), until
        a move fails or it has moved \emph{spread} cells; if it moved at least once, returns
        \code{true};
  \item otherwise (it has not moved) reverses its direction (\(\delta \leftarrow -\delta\)) and
        flows once more as in step~2; returns whether it moved.
\end{enumerate}

\task{D1} Write \code{Liquid}, \code{Water} and \code{Oil}, and add \verb|~| and \verb|o| to the
parser. Step 1 of rule L1 must be a call to \code{super.update(\ldots)} (R5): this is the
specialisation pattern of lesson~12.

\task{D2} Redefine \code{toString()} in \code{Liquid} so that it appends
\verb*| spread=|\textit{spread}\verb*| dir=|\textit{direction} to the general text, e.g.\
\verb|water '~' d=6 spread=3 dir=1|.

\paragraph{Checkpoint D.} Scene \code{layers} at ticks \(1\), \(2\), \(3\) (then nothing changes):
% scene layers
\begin{center}
\begin{checkpoint}{$t=1$}
\begin{BVerbatim}
#########
#       #
#.......#
#~~~~~~~#
#ooooooo#
#########
\end{BVerbatim}
\end{checkpoint}
\hspace{1.6em}
\begin{checkpoint}{$t=2$}
\begin{BVerbatim}
#########
#       #
#ooooooo#
#.......#
#~~~~~~~#
#########
\end{BVerbatim}
\end{checkpoint}
\hspace{1.6em}
\begin{checkpoint}{$t=3$}
\begin{BVerbatim}
#########
#       #
#ooooooo#
#~~~~~~~#
#.......#
#########
\end{BVerbatim}
\end{checkpoint}
\end{center}

Scene \code{dam} at ticks \(1\), \(2\), \(3\) and \(12\):
% scene dam
\begin{center}
\begin{checkpoint}{$t=1$}
\begin{BVerbatim}
############
#~  #      #
#~~~#      #
#~~~#      #
#~~~~ ~    #
############
\end{BVerbatim}
\end{checkpoint}
\hspace{1.6em}
\begin{checkpoint}{$t=2$}
\begin{BVerbatim}
############
#  ~#      #
#~~ #      #
#~~~#      #
#~~~~~   ~ #
############
\end{BVerbatim}
\end{checkpoint}
\hspace{1.6em}
\begin{checkpoint}{$t=3$}
\begin{BVerbatim}
############
#   #      #
#~~~#      #
#~~~#      #
#~~~~   ~ ~#
############
\end{BVerbatim}
\end{checkpoint}
\hspace{1.6em}
\begin{checkpoint}{$t=12$}
\begin{BVerbatim}
############
#   #      #
#   #      #
#~ ~#      #
#~~~~~~~~~~#
############
\end{BVerbatim}
\end{checkpoint}
\end{center}
After one tick of \code{dam}:
\begin{itemize}[nosep]
  \item \code{w.get(1, 1)} prints \verb|water '~' d=6 spread=3 dir=-1|.
\end{itemize}



\Q Nowhere in your code is there a line saying ``oil goes up''. Which rule makes oil end up above
the water in \code{layers}? Trace what happens to one oil cell and the water cell above it at
tick \(2\).

\Q Is the scene \code{dam} ever static? Argue from rule L1, then check it with \code{render()}.

\Q (\emph{Masquage in your own code.}) A colleague adds \code{private int direction = -1;} to
\code{Water}, ``so that water starts by flowing to the left''. Does the behaviour of water change?
What does \code{toString()} print for a fresh \code{Water}? Explain using A3.

% =====================================================================
\section{Part E --- Fire and smoke (45 min)}

\paragraph{Rule K1 (flammability).} \code{Wood} and \code{Oil} are flammable, with burn times
\(10\) and \(3\); \code{Water} extinguishes fire. Every other element keeps the general answers
of \code{Element} (\code{false}, \(0\), \code{false}).

\paragraph{Rule S1 (gases).} A \code{Gas} is a fluid particle with gravity \(-1\), a
\emph{lifetime} (given to its constructor) and an \emph{age}, initially \(0\). Its \code{update}:
first increments its age; if the age has reached the lifetime, it replaces itself by a new
\code{Air} (with \code{set}) and returns \code{true}; otherwise it performs the general movement of
rule P2. Its \code{toString()} appends \verb*| age=|\textit{age}\verb|/|\textit{lifetime}.

\paragraph{Rule K2 (fire).} A \code{Fire} has a number of remaining ticks. \code{Fire(int remaining)}
is the main constructor; \code{Fire()} uses \(6\) and must be written with \code{this(\ldots)}.
The orthogonal neighbours of \((r,c)\) are, in this order, up \((r-1,c)\), right \((r,c+1)\),
down \((r+1,c)\) and left \((r,c-1)\), ignoring those outside the grid. Its \code{update}:
\begin{enumerate}[nosep]
  \item if some neighbour extinguishes fire, the fire replaces itself by a new \code{Smoke} and
        returns \code{true};
  \item otherwise, each flammable neighbour \(n\), in the order above, is replaced by a new fire
        whose remaining ticks equal the burn time of \(n\);
  \item then the remaining ticks decrease by one; if they reach \(0\), the fire replaces itself by
        a new \code{Smoke} and returns \code{true}. Otherwise it returns \code{false}.
\end{enumerate}
Its \code{toString()} appends \verb*| left=|\textit{remaining}.

\task{E1} Write \code{Wood}, then add the redefinitions needed by rule K1. Remember R3 and R4:
\code{Fire} must work for any flammable element added later without being modified.

\task{E2} Write \code{Gas}, \code{Smoke} and \code{Fire}, and extend the parser with
\verb|=|, \verb|*| and \verb|"|.

\paragraph{Checkpoint E.} Scene \code{cabin} at ticks \(1\), \(2\), \(6\):
% scene cabin
\begin{center}
\begin{checkpoint}{$t=1$}
\begin{BVerbatim}
###############
#             #
#             #
#   =======   #
#   =     =   #
#   =  *  =   #
#   ===*===   #
###############
\end{BVerbatim}
\end{checkpoint}
\hspace{1.6em}
\begin{checkpoint}{$t=2$}
\begin{BVerbatim}
###############
#             #
#             #
#   =======   #
#   =     =   #
#   =  *  =   #
#   ==***==   #
###############
\end{BVerbatim}
\end{checkpoint}
\hspace{1.6em}
\begin{checkpoint}{$t=6$}
\begin{BVerbatim}
###############
#             #
#             #
#   =======   #
#   *     *   #
#   *  "  *   #
#   *******   #
###############
\end{BVerbatim}
\end{checkpoint}
\end{center}

then at ticks \(10\), \(15\), \(20\); everything has disappeared by tick \(30\):
% scene cabin
\begin{center}
\begin{checkpoint}{$t=10$}
\begin{BVerbatim}
###############
#             #
#             #
#   *******   #
#   *  "  *   #
#   *     *   #
#   *******   #
###############
\end{BVerbatim}
\end{checkpoint}
\hspace{1.6em}
\begin{checkpoint}{$t=15$}
\begin{BVerbatim}
###############
#             #
#             #
#   *******   #
#   *"" ""*   #
#  ""    """  #
#             #
###############
\end{BVerbatim}
\end{checkpoint}
\hspace{1.6em}
\begin{checkpoint}{$t=20$}
\begin{BVerbatim}
###############
#  """   """""#
#  "" " "  "  #
#    " "      #
#             #
#             #
#             #
###############
\end{BVerbatim}
\end{checkpoint}
\end{center}

Scene \code{oilfire} at ticks \(1\), \(2\), \(3\), \(6\):
% scene oilfire
\begin{center}
\begin{checkpoint}{$t=1$}
\begin{BVerbatim}
###########
#         #
#*        #
#*oooooooo#
#~~~~~~~~~#
###########
\end{BVerbatim}
\end{checkpoint}
\hspace{1.6em}
\begin{checkpoint}{$t=2$}
\begin{BVerbatim}
###########
#         #
#*        #
#*ooooooo"#
#~~~~~~~~~#
###########
\end{BVerbatim}
\end{checkpoint}
\hspace{1.6em}
\begin{checkpoint}{$t=3$}
\begin{BVerbatim}
###########
#         #
#*       "#
#*oooooo" #
#~~~~~~~~~#
###########
\end{BVerbatim}
\end{checkpoint}
\hspace{1.6em}
\begin{checkpoint}{$t=6$}
\begin{BVerbatim}
###########
#      """#
#"    "   #
#*ooo"    #
#~~~~~~~~~#
###########
\end{BVerbatim}
\end{checkpoint}
\end{center}

After one tick of \code{oilfire}:
\begin{itemize}[nosep]
  \item \code{w.get(2, 1)} prints \verb|fire '*' d=0 left=5|;
  \item \code{w.get(3, 1)} prints \verb|fire '*' d=0 left=3|;
  \item \code{w.get(3, 5)} prints \verb|oil 'o' d=4 spread=1 dir=-1|.
\end{itemize}
Explain why that oil cell has already turned around although it has not moved.

\Q \code{Particle.update} calls \code{getGravity()}, and was written before \code{Gas} existed.
When a smoke cell is updated, list the method bodies that execute (with their classes), and
explain why \code{Particle} does not need to be modified for smoke to rise.

\Q Suppose \code{set} did \emph{not} mark the new element (rule W1). Give a scene of one row and
a scene of one column in which a whole line of wood would burn in a single tick. Why does a fire
placed at the \emph{right} end of a row of wood not show the problem?

\Q Fire is hot and rises: would \code{Fire extends Gas} be a good idea? List what \code{Fire}
would inherit, what its \code{update} would have to look like, and whether the is-a test passes.

\Q \code{Wood} and \code{Oil} both redefine the flammability methods, yet their only common
super-class is \code{Element}. Is this duplication? Can it be removed with inheritance alone,
without making \code{Element} know about burning? (Keep your answer: interfaces come in week~5.)

% =====================================================================
\section{Part F --- Parallel futures: copying a world (30 min)}

Goal: from a world \code{w}, build an independent copy \code{future}, run the copy ahead, and
compare it with the present. A correct copy satisfies the following property, which is your test:
\[
  \textit{for every scene and all } k, n:\quad
  \text{run } k \text{ ticks, copy, run both } n \text{ ticks} \;\Longrightarrow\;
  \text{identical renders.}
\]

\task{F1} (\emph{Naive version.}) Write \code{World(World other)} so that it creates a
\emph{new} array but fills it with \emph{the same element references} as \code{other}. Run the
experiment: scene \code{cabin}, \(k=2\), \(n=5\). Then run it again with \code{hourglass},
\(k=2\), \(n=9\).

\Q One experiment passes and the other fails. Explain exactly which attributes are shared
between the two worlds, which method calls modify them, and why sand ``gets away with it''.
What does this tell you about the limits of testing?

\task{F2} (\emph{Deep version.}) Give every class a copy constructor \code{C(C other)} that
starts with \code{super(other)} and copies its own attributes (all state, including the mark).
Every class that can appear in a scene redefines
\begin{lstlisting}
public Element copy()   // new object of the same class, via the copy constructor
\end{lstlisting}
and \code{Element.copy()} returns \code{new Element(this)}. Rewrite \code{World(World other)} with
\code{copy()}. All experiments must now pass.

\Q In the copy constructor \code{Water(Water other)}, which calls \code{super(other)}, the argument has type \code{Water} but the
constructor of \code{Liquid} expects a \code{Liquid}. Why is this legal? Could the call
\code{new Element(other.get(r, c))} replace \code{copy()} in \code{World(World other)}? What would
the copied world look like after a few ticks?

\Q (\emph{What changes if\ldots?}) The copy constructor of \code{World} forgets to copy the tick.
In \code{hourglass} with \(n=9\), for which values of \(k\) does the experiment detect it, and what
does the difference look like? Predict first, then check with \(k=2\) and \(k=3\).

\Q Suppose \code{Water} forgets to redefine \code{copy()}. What is the dynamic type of a copied
water cell? (It depends on which of \code{Liquid}, \code{Particle} redefine \code{copy()}: answer
for your code and for the other possibilities.) With your code, which scenes of
Appendix~\ref{app:scenes} can detect the bug with the experiment, and which cannot? Under which
condition would the same mistake in \code{Sand} go unnoticed in \emph{every} possible scene?

% =====================================================================
\section{Part G --- Your scene, the design table, reflection (20 min)}

\task{G1} Using the helper of Appendix~\ref{app:helper}, write a \code{main} that animates a scene:
clear the screen, print \code{render()} and the tick, pause about 80\,ms, step, repeat.

\task{G2} Design your own scene, at least \(30 \times 15\), using at least six kinds of element,
that tells a small story in at most 200 ticks (for instance, an oil tank leaks, catches fire, and a
water tower bursts in time, or not). Keep it in your main class as a \code{String[]}.

\task{G3} Fill in a table with one row per class and one column per method of the
\code{Element} interface (plus \code{getGravity}, \code{canEnter}, \code{tryMove}, \code{copy}).
In each cell write D (defined here), R (redefined here) or nothing (inherited). Then answer:

\Q Count the R entries for \code{update}. For each, say whether the redefinition \emph{replaces}
or \emph{extends} (with \code{super}) the general behaviour.

\Q \code{Element.update} has an empty body and \verb|new Element("?", '?', 0)| is accepted by the
compiler, although such an object is meaningless. Describe the two problems this causes, and the
language feature you would like to fix them.

\Q All your classes live in the same (default) package. In this setting, what does
\code{protected} give compared with \code{private}, and compared with \code{public}? Would making
the attributes of \code{Element} \code{protected} have saved you code? Would it have been worth it?

% =====================================================================
\section*{Optional extensions}

Each extension is a test of your design: count how many existing classes you had to modify.
\begin{itemize}
  \item \textbf{Steam.} When water extinguishes a fire, the fire becomes \code{Steam} (a gas)
        instead of smoke, and steam turns back into water when its lifetime ends. Rule S1 creates
        \code{Air} at the end of a gas's life: refactor \code{Gas} so that a sub-class can choose
        what it leaves behind, by redefining a single method.
  \item \textbf{Lava.} A liquid, as dense as sand, that ignites its flammable neighbours like a
        fire and turns into stone when it touches water. Which code of \code{Fire} would you like
        to reuse, and why does single inheritance prevent it?
  \item \textbf{Salt.} A powder that disappears when it is next to water, which becomes
        \code{SaltWater}: denser than water, so it sinks below it.
\end{itemize}

\appendix
% =====================================================================
\section{Scenes}\label{app:scenes}

Copy each array into your main class. PDF viewers usually drop the spaces inside these strings
when you copy them, so copy the arrays from the \LaTeX{} source of this document instead.

\noindent
\begin{minipage}[t]{0.32\linewidth}
\begin{lstlisting}
String[] drop = {
    "#######",
    "#  .  #",
    "#  .  #",
    "#  .  #",
    "#     #",
    "#     #",
    "#######"
};
\end{lstlisting}
\end{minipage}\hfill
\begin{minipage}[t]{0.32\linewidth}
\begin{lstlisting}
String[] hourglass = {
    "###########",
    "#.........#",
    "##.......##",
    "###.....###",
    "####...####",
    "##### #####",
    "####   ####",
    "###     ###",
    "##       ##",
    "#         #",
    "###########"
};
\end{lstlisting}
\end{minipage}\hfill
\begin{minipage}[t]{0.32\linewidth}
\begin{lstlisting}
String[] layers = {
    "#########",
    "#.......#",
    "#~~~~~~~#",
    "#ooooooo#",
    "#       #",
    "#########"
};
\end{lstlisting}
\end{minipage}

\noindent
\begin{minipage}[t]{0.32\linewidth}
\begin{lstlisting}
String[] dam = {
    "############",
    "#~~~#      #",
    "#~~~#      #",
    "#~~~#      #",
    "#~~~       #",
    "############"
};
\end{lstlisting}
\end{minipage}\hfill
\begin{minipage}[t]{0.32\linewidth}
\begin{lstlisting}
String[] cabin = {
    "###############",
    "#             #",
    "#             #",
    "#   =======   #",
    "#   =     =   #",
    "#   =  *  =   #",
    "#   =======   #",
    "###############"
};
\end{lstlisting}
\end{minipage}\hfill
\begin{minipage}[t]{0.32\linewidth}
\begin{lstlisting}
String[] oilfire = {
    "###########",
    "#         #",
    "#*        #",
    "#ooooooooo#",
    "#~~~~~~~~~#",
    "###########"
};
\end{lstlisting}
\end{minipage}

\section{Terminal helper}\label{app:helper}

Paste these two functions into your main class and treat them as a black box: \code{try}/\code{catch}
belongs to week~6. If your terminal does not understand the escape sequence (some Windows
consoles), print about forty empty lines instead.
\begin{lstlisting}
static void clearScreen() {
    System.out.print("\033[H\033[2J");
    System.out.flush();
}

static void pause(int milliseconds) {
    try { Thread.sleep(milliseconds); } catch (InterruptedException e) { }
}
\end{lstlisting}

\section{Self-assessment}

\begin{center}
\small
\begin{tabular}{@{}>{\raggedright}p{5.4cm}l>{\raggedright\arraybackslash}p{6.4cm}@{}}
\toprule
Concept & Lesson & Where it is tested \\
\midrule
\code{extends}, is-a vs has-a, transitivity & 10 & Q9, Q12, Q18, the hierarchy itself \\
\code{protected}, packages & 11 & Q7, Q8, Q26, R2 \\
redefinition and \code{super.}\emph{method} & 12 & Q3, Q4, Tasks D1, D2, E1, Q24 \\
\emph{masquage} & 12 & Q7, Q15 \\
\code{super(\ldots)}, \code{this(\ldots)}, order of construction & 13 & Q1, Q2, Q6, Q10, R6, rule K2, F2 \\
static vs dynamic binding & 14 & Q3--Q5, Q7, Q16, Q21, Q23 \\
copy constructors, references & week 2 & F1, F2, Q20--Q23 \\
2-D arrays, strings, loops, functions & BOOC I & B3, C2, G1, G2 \\
\bottomrule
\end{tabular}
\end{center}

\end{document}